\begin{table*}[!ht]
\centering
\caption*{\textbf{Supplementary Table~S1. The 39 second-level (L2) subtypes of the
single-nucleus atlas.} Subtypes are listed in descending order of cell number, continued
across the two halves of the table. Counts refer to the 399{,}579 nuclei retained after
quality control and are the reference used for spatial deconvolution. These names are not
unique identifiers of a lineage--subtype combination: they map onto 67 distinct
combinations, and \emph{Basophil/Mast~1} is assigned predominantly to epithelial rather
than myeloid cells (main text).}
\label{tab:s1}
\small
\begin{tabular}{@{}lr@{\hspace{3.2em}}lr@{}}
\toprule
\tabheadfont Subtype & \tabheadfont Nuclei & \tabheadfont Subtype & \tabheadfont Nuclei \\
\midrule
AT2 & 88,666 & Vein & 5,654 \\
CD4$^{+}$ Naive T & 37,528 & Bronchial Vessel 2 & 4,674 \\
Fibroblast & 35,581 & Goblet/Mucous & 3,587 \\
AT1 & 30,320 & Pericyte & 3,480 \\
Myofibroblast & 24,244 & Lymphatic & 3,354 \\
TREM2$^{+}$ Dendritic & 20,571 & CD8$^{+}$ Naive T & 2,680 \\
Plasma & 15,975 & EREG$^{+}$ Dendritic & 2,476 \\
Macrophage & 10,837 & Basal & 2,386 \\
NK & 10,608 & Nonclassical Monocyte & 2,327 \\
IGSF21$^{+}$ Dendritic & 9,193 & Bronchial Vessel 1 & 2,292 \\
Capillary & 8,839 & NKT & 2,172 \\
Vascular Smooth Muscle & 8,408 & Classical Monocyte & 1,985 \\
CD8$^{+}$ Mem/Eff T & 8,173 & Myeloid Dendritic Type 1 & 1,953 \\
Capillary Intermediate 1 & 7,194 & CD4$^{+}$ Mem/Eff & 1,917 \\
B & 6,940 & Serous & 1,728 \\
Ciliated & 6,697 & Plasmacytoid Dendritic & 925 \\
Artery & 6,388 & Neutrophil & 840 \\
Basophil/Mast 1 & 6,368 & Capillary Aerocyte & 305 \\
Myeloid Dendritic Type 2 & 6,100 & Proliferating Macrophage & 294 \\
Lipofibroblast & 5,920 &  &  \\
\bottomrule
\end{tabular}
\end{table*}
